\documentclass{article}
\usepackage[T1]{fontenc}
\usepackage{lmodern}
\usepackage{graphicx}
\usepackage{amsmath,amssymb}
\usepackage[a4paper,margin=2cm]{geometry}
\usepackage[absolute,overlay]{textpos}
\setlength{\TPHorizModule}{1mm}
\setlength{\TPVertModule}{1mm}
\usepackage[indonesian]{babel}
\usepackage{hyperref}
\setlength{\emergencystretch}{2em}
% unit: Halaman 29

\begin{document}

% === Halaman 29 (llm) ===

\section*{Kode Program RSA dalam Bahasa Python}

\begin{enumerate}
    \item Kode program enkripsi dan dekripsi pesan dengan kunci publik dan kunci privat milik sendiri, di dalamnya termasuk:
    \begin{itemize}
        \item mensimulasikan perhitungan kunci publik dan kunci privat
    \end{itemize}
    \item Kode program enkripsi dan dekripsi pesan dengan kunci publik dan kunci privat milik sendiri, di dalamnya termasuk:
    \begin{itemize}
        \item membangkitkan kunci publik dan kunci privat dengan fungsi RSA.generate()
        \item Kunci publiK dan kunci privat disimpan di dalam file dengan format PEM
    \end{itemize}
    \item Kode program enkripsi pesan dengan kunci publik pengirim (Bob) dan dekripsi pesan dengan kunci privat penerima pesan (Alice)
\end{enumerate}

\end{document}
